% !TEX root = ../main.tex

\chapter{توصیف و طراحی سامانه}

% =========================================================
% 4-1
% =========================================================
\section{مقدمه}

در این فصل، معماری کلی سامانه گفت‌وگوی صوتی فارسی طراحی‌شده
در این پروژه تشریح می‌شود.

هدف اصلی سامانه، دریافت گفتار فارسی از کاربر، تبدیل آن به متن،
پردازش متن توسط یک مدل زبانی و در نهایت تولید پاسخ صوتی فارسی
است.

معماری پایه سامانه به‌صورت یک خط لوله ماژولار طراحی شده است:

\begin{center}
\lr{Speech}
$\longrightarrow$
\lr{ASR}
$\longrightarrow$
\lr{LLM}
$\longrightarrow$
\lr{TTS}
$\longrightarrow$
\lr{Speech}
\end{center}

علاوه بر این سه مؤلفه اصلی، ماژول‌هایی مانند تشخیص فعالیت
صوتی، پیش‌پردازش صوت، مدیریت تاریخچه مکالمه، پردازش متن و
مدیریت خطا نیز در طراحی سامانه در نظر گرفته شده‌اند.

یکی از اصول اصلی طراحی، قابلیت جایگزینی مستقل هر مدل است.
به این ترتیب، می‌توان بدون تغییر اساسی در کل سامانه، مدل \lr{ASR}،
\lr{LLM} یا \lr{TTS} را تغییر داد.

در این فصل ابتدا نیازمندی‌های طراحی معرفی می‌شوند و سپس معماری
کلی، جریان داده، رابط ماژول‌ها و تصمیم‌های اصلی طراحی تشریح
خواهند شد.


% =========================================================
% 4-2
% =========================================================
\section{نیازمندی‌های سامانه}

نیازمندی‌های پروژه را می‌توان به دو گروه اصلی تقسیم کرد:
نیازمندی‌های عملکردی و نیازمندی‌های غیرعملکردی.


% ---------------------------------------------------------
\subsection{نیازمندی‌های عملکردی}

سامانه باید بتواند وظایف اصلی زیر را انجام دهد:

\begin{itemize}

    \item دریافت صوت فارسی از کاربر؛

    \item تشخیص محتوای گفتار و تبدیل آن به متن؛

    \item ارسال متن به مدل زبانی؛

    \item تولید پاسخ متنی فارسی؛

    \item تبدیل پاسخ متنی به صوت؛

    \item پخش پاسخ صوتی برای کاربر؛

    \item و نگهداری اطلاعات مورد نیاز از تاریخچه مکالمه.

\end{itemize}

همچنین طراحی باید امکان افزودن مؤلفه‌هایی مانند
\lr{VAD}
و کاهش نویز را فراهم کند.


% ---------------------------------------------------------
\subsection{نیازمندی‌های غیرعملکردی}

در کنار عملکرد صحیح، چند ویژگی غیرعملکردی نیز برای سامانه
اهمیت دارند.

نخست، سامانه باید تا حد امکان روی سخت‌افزار شخصی و بدون
وابستگی ضروری به سرویس‌های ابری اجرا شود.

دوم، تأخیر باید به اندازه‌ای محدود باشد که مکالمه برای کاربر
قابل قبول باقی بماند.

سوم، مصرف منابع باید کنترل شود، زیرا چند مدل هوش مصنوعی
ممکن است هم‌زمان در حافظه قرار گیرند.

چهارم، طراحی باید ماژولار باشد تا ارزیابی و جایگزینی مدل‌ها
ساده باشد.

در نتیجه، مهم‌ترین نیازمندی‌های غیرعملکردی عبارت‌اند از:

\begin{itemize}

    \item اجرای محلی؛
    \item تأخیر پایین؛
    \item مصرف مناسب حافظه؛
    \item ماژولار بودن؛
    \item قابلیت توسعه؛
    \item و امکان ارزیابی مستقل هر مؤلفه.
\end{itemize}


% =========================================================
% 4-3
% =========================================================
\section{معماری کلی سامانه}

معماری سامانه به‌صورت زنجیره‌ای و ماژولار طراحی شده است.

در سطح مفهومی، جریان اصلی پردازش به شکل زیر است:

\begin{center}
Microphone
$\longrightarrow$
\lr{Audio} Processing
$\longrightarrow$
\lr{ASR}
$\longrightarrow$
\lr{Text} Processing
$\longrightarrow$
Conversation Manager
$\longrightarrow$
\lr{LLM}
$\longrightarrow$
\lr{TTS}
$\longrightarrow$
Speaker
\end{center}

هر بخش دارای ورودی و خروجی مشخصی است و تنها از طریق رابط
تعریف‌شده با بخش بعدی ارتباط برقرار می‌کند.

این طراحی باعث می‌شود وابستگی مستقیم میان پیاده‌سازی داخلی
مدل‌ها کاهش پیدا کند.


% ---------------------------------------------------------
\subsection{ماژول ورودی صوت}

ماژول ورودی صوت مسئول دریافت صدای کاربر از میکروفون است.

خروجی این بخش یک جریان یا فایل صوتی با قالب مشخص است که
برای پردازش‌های بعدی استفاده می‌شود.

در طراحی نهایی، این بخش می‌تواند علاوه بر ضبط صوت، وظایفی
مانند:

\begin{itemize}

    \item تنظیم نرخ نمونه‌برداری؛
    \item تبدیل تعداد کانال‌ها؛
    \item نرمال‌سازی سطح صوت؛
    \item و آماده‌سازی فرمت مناسب برای \lr{ASR}
\end{itemize}

را نیز انجام دهد.


% ---------------------------------------------------------
\subsection{ماژول VAD و پیش‌پردازش صوت}

ماژول
\lr{VAD}
وظیفه تشخیص وجود گفتار در جریان صوتی را بر عهده دارد.

هدف این ماژول جلوگیری از ارسال بخش‌های طولانی سکوت یا
نویز غیرضروری به مدل \lr{ASR} است.

به‌صورت کلی:

\begin{center}
\lr{Audio} Stream
$\longrightarrow$
VAD
$\longrightarrow$
\lr{Speech} Segment
\end{center}

در کنار VAD، امکان افزودن یک ماژول کاهش نویز نیز در طراحی
در نظر گرفته شده است.

این بخش به‌صورت اختیاری طراحی شده تا بتوان اثر آن را بر
دقت \lr{ASR} و تأخیر کلی سامانه جداگانه ارزیابی کرد.


% ---------------------------------------------------------
\subsection{ماژول \lr{ASR}}

ماژول \lr{ASR} صوت ورودی را دریافت کرده و متن فارسی تولید می‌کند.

رابط مفهومی این بخش به شکل زیر است:

\begin{center}
\lr{Audio}
$\longrightarrow$
\lr{ASR} \lr{Model}
$\longrightarrow$
Transcription
\end{center}

از آنجا که چند مدل \lr{ASR} در پروژه ارزیابی شده‌اند، این ماژول
به‌گونه‌ای طراحی می‌شود که مدل داخلی قابل جایگزینی باشد.

بنابراین بخش‌های دیگر سامانه نباید به جزئیات معماری یک مدل
خاص وابسته باشند.


% ---------------------------------------------------------
\subsection{ماژول پردازش متن}

متن تولیدشده توسط \lr{ASR} ممکن است پیش از ارسال به مدل زبانی
نیازمند پردازش باشد.

این ماژول می‌تواند وظایفی مانند:

\begin{itemize}

    \item نرمال‌سازی نویسه‌ها؛
    \item مدیریت فاصله و نیم‌فاصله؛
    \item حذف خروجی‌های خالی یا نامعتبر؛
    \item مدیریت علائم نگارشی؛
    \item و آماده‌سازی متن برای مدل زبانی
\end{itemize}

را انجام دهد.

این بخش میان \lr{ASR} و مدل زبانی قرار می‌گیرد و هدف آن کاهش
وابستگی مدل زبانی به قالب خروجی یک \lr{ASR} خاص است.


% ---------------------------------------------------------
\subsection{مدیریت مکالمه}

برای پشتیبانی از مکالمه چندمرحله‌ای، سامانه نیازمند یک مؤلفه
برای نگهداری تاریخچه گفتگو است.

در ساده‌ترین حالت، تاریخچه شامل پیام‌های قبلی کاربر و پاسخ‌های
مدل است:

\begin{equation}
H_t =
\{U_1,A_1,U_2,A_2,\ldots,U_t\}
\end{equation}

که در آن $U_i$ پیام کاربر و $A_i$ پاسخ سامانه است.

هنگام ارسال درخواست جدید، بخشی از این تاریخچه همراه با پیام
جدید در اختیار مدل زبانی قرار می‌گیرد.

به دلیل محدودیت طول بافت مدل، تاریخچه نمی‌تواند بدون محدودیت
رشد کند.
بنابراین طراحی مدیریت مکالمه باید امکان حذف، محدودکردن یا
خلاصه‌سازی پیام‌های قدیمی را فراهم کند.


% ---------------------------------------------------------
\subsection{ماژول مدل زبانی}

ماژول مدل زبانی متن پردازش‌شده و تاریخچه مکالمه را دریافت کرده
و پاسخ متنی تولید می‌کند.

رابط آن به‌صورت مفهومی چنین است:

\begin{center}
Prompt + Context
$\longrightarrow$
\lr{LLM}
$\longrightarrow$
Response \lr{Text}
\end{center}

مانند \lr{ASR}، این بخش نیز به‌صورت قابل جایگزینی طراحی شده است
تا بتوان مدل‌هایی با اندازه و نیاز سخت‌افزاری متفاوت را
آزمایش کرد.


% ---------------------------------------------------------
\subsection{ماژول \lr{TTS}}

ماژول \lr{TTS} پاسخ مدل زبانی را دریافت کرده و آن را به صوت فارسی
تبدیل می‌کند.

ساختار مفهومی آن عبارت است از:

\begin{center}
Response \lr{Text}
$\longrightarrow$
\lr{TTS} \lr{Model}
$\longrightarrow$
\lr{Audio} Response
\end{center}

مدل \lr{TTS} نیز باید بدون تغییر سایر بخش‌های سامانه قابل تعویض باشد.

در این مرحله ممکن است پیش از تولید صوت، نرمال‌سازی متن
مخصوص \lr{TTS} نیز انجام شود؛ برای مثال تبدیل عدد یا علامت به
نمایش مناسب برای تلفظ.


% =========================================================
% 4-4
% =========================================================
\section{جریان داده در سامانه}

فرآیند کامل یک نوبت مکالمه را می‌توان در چند مرحله خلاصه کرد.

ابتدا کاربر صحبت می‌کند و صوت توسط میکروفون دریافت می‌شود.
پس از تشخیص بخش گفتاری، صوت برای مدل \lr{ASR} ارسال می‌شود.

خروجی متنی \lr{ASR} پس از نرمال‌سازی همراه با تاریخچه مکالمه
برای مدل زبانی آماده می‌شود.

مدل زبانی پاسخ متنی را تولید می‌کند و این پاسخ ابتدا در
تاریخچه مکالمه ذخیره شده و سپس به مدل \lr{TTS} ارسال می‌شود.

در نهایت، صوت تولیدشده برای کاربر پخش می‌شود.

به شکل خلاصه:

\begin{center}
User
$\rightarrow$
\lr{Audio}
$\rightarrow$
\lr{ASR}
$\rightarrow$
\lr{Text}
$\rightarrow$
Context
$\rightarrow$
\lr{LLM}
$\rightarrow$
\lr{Text}
$\rightarrow$
\lr{TTS}
$\rightarrow$
\lr{Audio}
$\rightarrow$
User
\end{center}